\chapter*{Glossary}
\addcontentsline{toc}{chapter}{Glossary}
These entries identify distinct interfaces used in the current manuscript.
A theorem's full domain and assumptions remain in its owning chapter.
\begin{description}
\item[Sequence convention] Unless an alignment explicitly labels opposite
polarity, a strand's sequence is written in its own $5'$ to $3'$ direction.
\item[Reverse complement] Complement each base and reverse order to retain that
same sequence convention. An aligned complementary strand is a different display.
\item[Domain] A contiguous segment assigned a recognition role in a coarse
molecular design, not a universal rate constant or a proof of specificity.
\item[Hybridization] Association through complementary interactions; pairing
does not seal a covalent backbone nick.
\item[Ligation] Covalent joining of suitable adjacent strand ends. Required
chemistry and material endpoints differ from hybridization.
\item[Support] Identities with nonzero population multiplicity.
\item[Multiplicity] Number of copies of an identity, distinct from the number
of distinct identities.
\item[Certificate] A proposed witness checked against the original instance.
\item[Soundness] Every accepted witness satisfies the declared problem and decoder.
\item[Completeness] Required existing witnesses can be recovered under the
declared algorithm; physical sampling and detection add separate failure events.
\item[Reaction extent] Amount of net conversion in declared concentration units.
\item[Stoichiometric matrix] Product minus reactant coefficients; it defines
net change but does not alone specify reaction enabledness.
\item[Propensity] Instantaneous reaction-event hazard for a discrete count state.
\item[Linear invariant] A weighted inventory annihilating the stoichiometric
matrix, hence preserved by each permitted reaction firing.
\item[First passage] First arrival at a specified outcome set; successful and
failed attempts can have different conditional durations.
\item[Stable output] An output that cannot change under permitted continuations;
this need not mean the entire reaction network has halted.
\item[Simulation invariant] Relation between source and target representations
preserved by each simulated transition.
\item[Uniform construction] Effective generation of a device family under a
declared construction and preparation cost model.
\item[Resource coupling] Dependence caused by a shared finite inventory, even
without unintended sequence recognition.
\item[Directed assembly] In the stated tile model, every producible assembly
extends toward the same terminal assembly; physical yield needs additional evidence.
\item[Common-mode fault] A shared failure mechanism producing dependent errors;
independent-vote formulas need not apply.
\item[Poissonized preparation] Independent random copy counts under the declared
model, with a random total count rather than a fixed draw total.
\item[Coverage] All required identities are represented; abundant total DNA
alone does not certify this event.
\item[Readout contract] Sampling, detection, calibration, decoder and output
interface specifying how physical state becomes a reported answer.
\item[Pareto dominance] No declared controlled cost is worse and at least one
is better under matched correctness, deadline and observation requirements.
\item[Located erasure] An unavailable symbol or shard whose position/address is
known; this differs from an unlocated erroneous value.
\item[Trusted manifest] Out-of-band source identity, length and digest against
which reconstruction is checked; a checksum alone does not authenticate it.
\item[Enumerative code] Reversible mapping between legal words and lexical ranks
using exact suffix-completion counts under declared finite constraints.
\item[Logical storage rate] Source bits divided by unique encoded nucleotides;
physical-copy rate counts all molecule copies in its denominator.
\item[Molecular address] Recognition/selection mechanism determining which
molecules are observed, distinct from a digital index interpreted after readout.
\item[Evidence ledger] Byte-bound, ordered provenance records; structural validity
does not authenticate an author or establish experimental truth.
\item[Preparation identity] Identity of a specific physical preparation, distinct
from an abstract design, sequence realization, material lot or observation.
\item[Independent experimental unit] Unit independently assigned or prepared under
the declared design; nested observations do not automatically increase its count.
\item[Observation freeze] Recorded boundary fixing raw evidence and measurement
mapping before the declared unblinding or analytical changes.
\item[Experimental replication] New independently prepared physical evidence,
distinct from computational reproduction of a frozen analysis.
\item[Enabledness] Every consumed multiplicity must be present, including catalysts
with zero net stoichiometric increment.
\item[State equation] Necessary relation $x=x_0+S\sigma$ for rule firing counts;
it does not establish a legal firing order or reachable target.
\item[Search closure] A completely processed discovered set closed under every
enabled successor; distinct from a depth-bound or interrupted budget.
\item[Labeled transition] State/rule/state triple, preserved even when rules share
targets or return to the same state.
\item[Propensity] State-conditioned hazard in events per time for a declared reaction channel.
\item[Count trajectory] Right-continuous integer state path with event times; distinct from concentration or ensemble mean.
\item[Event-budget censoring] Simulation stopped before its horizon; its last state is not a completed horizon observation.
\item[Positivity bound] Step condition preserving nonnegative populations; distinct from linear absolute stability.
\item[Factorial moment] Expectation of count products such as $A(A-1)$; generally not a product evaluated at the mean.
\end{description}
